% GENERATED FILE. Edit canonical lessons, then run npm run generate:books.
% canonical-source-hash: fnv1a64:6ddf1719fcf601a7
% canonical-lessons: ES-C528-desastre, ES-C528-catastrofe, ES-C528-tragedia, ES-C528-terremoto, ES-C528-inundacion

\chapter{Una catástrofe}
\label{ch:es-una-catastrofe}
\hlchaptermodality{528}

\begin{chapteropening}
\textbf{By the end of this chapter:} \emph{I can talk about a disaster in the news: an earthquake, a flood, a tragedy.}

Complete the last lesson of chapter 528: I can talk about a disaster in the news: an earthquake, a flood, a tragedy.
\end{chapteropening}

\section[el desastre]{el desastre — the disaster}

\label{lesson:ES-C528-desastre}



\begin{quote}
Before the new one: say the Spanish for unhurt, then the Spanish for to fall.
\end{quote}

\subsection*{You'll want to know: el desastre}
\textbf{el desastre} — \textquotedblleft{}the disaster\textquotedblright{}.

\begin{quote}
La fiesta fue un desastre.
\end{quote}

\begin{quote}
\emph{The party was a disaster.}
\end{quote}

In everyday speech it is anything that goes badly: \emph{¡Qué desastre!}

\begin{cousinweb}
\emph{desastre} is usually traced to Occitan \emph{desastre} or Italian \emph{disastro}, where it meant 'bad star', built on \emph{des-} and the word for 'star', from Latin \emph{astrum} (Lewis \& Short). A disaster was blamed on the stars.
\end{cousinweb}

\begin{grammarlens}[title={What you've built}]
One more word for telling, explaining and describing at B1.
\end{grammarlens}

\subsection*{Your turn}
\begin{itemize}
\raggedright
  \item \emph{Say it:} \emph{el desastre}
  \item \emph{Say it:} \emph{el desastre}, once more
  \item \emph{Say it:} a sentence of your own with \emph{el desastre}
  \item \emph{Recall:} say the Spanish for unhurt, then the Spanish for to fall, then say \emph{el desastre} again
\end{itemize}

\begin{tcolorbox}[breakable,colback=teal!4,colframe=teal!35!black,title={Before you move on}]
What does el desastre mean? (\textquotedblleft{}The disaster\textquotedblright{}.) Say it once more.
\end{tcolorbox}
\section[la catástrofe]{la catástrofe — the catastrophe}

\label{lesson:ES-C528-catastrofe}



\begin{quote}
Before the new one: say the Spanish for to fall, then the Spanish for the disaster.
\end{quote}

\subsection*{You'll want to know: la catástrofe}
\textbf{la catástrofe} — \textquotedblleft{}the catastrophe\textquotedblright{}.

\begin{quote}
El huracán fue una catástrofe para la isla.
\end{quote}

\begin{quote}
\emph{The hurricane was a catastrophe for the island.}
\end{quote}

\begin{cousinweb}
\emph{catástrofe} is a learned borrowing that comes from Greek \emph{katastroph\'{\={e}}} 'an overturning' (EtymDB), built on \emph{katá} 'down' and \emph{stréphein} 'to turn'.
\end{cousinweb}

\begin{grammarlens}[title={What you've built}]
One more word for telling, explaining and describing at B1.
\end{grammarlens}

\subsection*{Your turn}
\begin{itemize}
\raggedright
  \item \emph{Say it:} \emph{la catástrofe}
  \item \emph{Say it:} \emph{la catástrofe}, once more
  \item \emph{Say it:} a sentence of your own with \emph{la catástrofe}
  \item \emph{Recall:} say the Spanish for to fall, then the Spanish for the disaster, then say \emph{la catástrofe} again
\end{itemize}

\begin{tcolorbox}[breakable,colback=teal!4,colframe=teal!35!black,title={Before you move on}]
What does la catástrofe mean? (\textquotedblleft{}The catastrophe\textquotedblright{}.) Say it once more.
\end{tcolorbox}
\section[la tragedia]{la tragedia — the tragedy}

\label{lesson:ES-C528-tragedia}



\begin{quote}
Before the new one: say the Spanish for the disaster, then the Spanish for the catastrophe.
\end{quote}

\subsection*{You'll want to know: la tragedia}
\textbf{la tragedia} — \textquotedblleft{}the tragedy\textquotedblright{}.

\begin{quote}
El accidente fue una tragedia para el pueblo.
\end{quote}

\begin{quote}
\emph{The accident was a tragedy for the town.}
\end{quote}

\begin{cousinweb}
\emph{tragedia} is a learned borrowing from Latin \emph{tragoedia}, which comes from Greek \emph{tragōidía}, built on \emph{trágos} 'goat' and \emph{ōid\'{\={e}}} 'song' (EtymDB).
\end{cousinweb}

\begin{grammarlens}[title={What you've built}]
One more word for telling, explaining and describing at B1.
\end{grammarlens}

\subsection*{Your turn}
\begin{itemize}
\raggedright
  \item \emph{Say it:} \emph{la tragedia}
  \item \emph{Say it:} \emph{la tragedia}, once more
  \item \emph{Say it:} a sentence of your own with \emph{la tragedia}
  \item \emph{Recall:} say the Spanish for the disaster, then the Spanish for the catastrophe, then say \emph{la tragedia} again
\end{itemize}

\begin{tcolorbox}[breakable,colback=teal!4,colframe=teal!35!black,title={Before you move on}]
What does la tragedia mean? (\textquotedblleft{}The tragedy\textquotedblright{}.) Say it once more.
\end{tcolorbox}
\section[el terremoto]{el terremoto — the earthquake}

\label{lesson:ES-C528-terremoto}



\begin{quote}
Before the new one: say the Spanish for the catastrophe, then the Spanish for the tragedy.
\end{quote}

\subsection*{You'll want to know: el terremoto}
\textbf{el terremoto} — \textquotedblleft{}the earthquake\textquotedblright{}.

\begin{quote}
El terremoto duró treinta segundos.
\end{quote}

\begin{quote}
\emph{The earthquake lasted thirty seconds.}
\end{quote}

\emph{Un sismo} is a more technical word, common in Latin America.

\begin{cousinweb}
\emph{terremoto} was borrowed from Italian \emph{terremoto}, from Latin \emph{terrae mōtus} 'a movement of the earth' (EtymDB), built on \emph{terra} 'earth' and \emph{mōtus} 'movement'.
\end{cousinweb}

\begin{grammarlens}[title={What you've built}]
One more word for telling, explaining and describing at B1.
\end{grammarlens}

\subsection*{Your turn}
\begin{itemize}
\raggedright
  \item \emph{Say it:} \emph{el terremoto}
  \item \emph{Say it:} \emph{el terremoto}, once more
  \item \emph{Say it:} a sentence of your own with \emph{el terremoto}
  \item \emph{Recall:} say the Spanish for the catastrophe, then the Spanish for the tragedy, then say \emph{el terremoto} again
\end{itemize}

\begin{tcolorbox}[breakable,colback=teal!4,colframe=teal!35!black,title={Before you move on}]
What does el terremoto mean? (\textquotedblleft{}The earthquake\textquotedblright{}.) Say it once more.
\end{tcolorbox}
\section[la inundación]{la inundación — the flood}

\label{lesson:ES-C528-inundacion}



\begin{quote}
Before the new one: say the Spanish for the tragedy, then the Spanish for the earthquake.
\end{quote}

\subsection*{You'll want to know: la inundación}
\textbf{la inundación} — \textquotedblleft{}the flood\textquotedblright{}.

\begin{quote}
Después de la lluvia hubo inundaciones.
\end{quote}

\begin{quote}
\emph{After the rain there were floods.}
\end{quote}

\begin{cousinweb}
\emph{inundación} is a learned borrowing from Latin \emph{inundātiō} 'an overflowing' (EtymDB), built on \emph{unda} 'a wave'.
\end{cousinweb}

\begin{grammarlens}[title={What you've built}]
One more word for telling, explaining and describing at B1.
\end{grammarlens}

\subsection*{Your turn}
\begin{itemize}
\raggedright
  \item \emph{Say it:} \emph{la inundación}
  \item \emph{Say it:} \emph{la inundación}, once more
  \item \emph{Say it:} a sentence of your own with \emph{la inundación}
  \item \emph{Recall:} say the Spanish for the tragedy, then the Spanish for the earthquake, then say \emph{la inundación} again
\end{itemize}

\begin{tcolorbox}[breakable,colback=teal!4,colframe=teal!35!black,title={Before you move on}]
What does la inundación mean? (\textquotedblleft{}The flood\textquotedblright{}.) Say it once more.
\end{tcolorbox}
